\chapter{Estado del Arte de Tecnologías para Desarrollo Web}
\label{cap:estadoDelArteDeTecnologias}


\section{Introducción}

Las aplicaciones web modernas se caracterizan por su dinamismo, interactividad, y capacidad de operar en tiempo real. Estas propiedades no son completamente nuevas, sino el resultado de una evolución progresiva del desarrollo web\footnote{En el siguiente artículo se puede leer más sobre el progreso histórico del desarrollo web: \url{https://www.geeksforgeeks.org/blogs/history-and-evolution-of-web-development/}}, en la que las páginas estáticas y dinámicas han coexistido y se han complementado desde finales del siglo XX. Mientras que las páginas estáticas siguen siendo adecuadas para ciertos casos de uso, las aplicaciones web actuales tienden a incorporar mecánismos más avanzados de interacción. 

En sus inicios, la web estaba concebida como un sistema de documentos estáticos interconectados mediante enlaces: el servidor enviaba archivos HTML que el cliente (navegador) simplemente renderizaba. En este sentido, el modelo original puede entenderse como una extensión del intercambio de documentos. Con el tiempo, este modelo básico se amplió mediante dos avances fundamentales. Por un lado, la generación dinámica de contenido en el servidor permitió construir HTML bajo demanda en función de las peticiones del usuario. Por otro lado, la ejecución de código en el lado del cliente introdujo la posibilidad de modificar la interfaz y el comportamiento de la aplicación sin necesidad de nuevas comunicaciones con el servidor. Estas dos mejoras sentaron las bases de la separación conceptual entre cliente y servidor, que posteriormente daría lugar a los términos \textit{frontend} y \textit{backend}, entendidos como los componentes que se ejecutan en cada uno de estos lados.

Para cumplir con las expectativas de los usuarios y las exigencias del mercado, las aplicaciones web deben satisfacer una serie de requisitos fundamentales. Entre ellos, destacan el alto rendimiento, la escalabilidad y la mantenibilidad del código. Aunque la seguridad también puede ser una prioridad crítica dependiendo del contexto de la aplicación, estos tres primeros requisitos son universales en el desarrollo web actual.

En lo que respecta al rendimiento, las aplicaciones deben operar con rapidez y eficiencia. En 2016, la mitad de los usuarios de la web esperaban que una página se cargase en menos de dos segundos\footnote{\url{https://blog.google/products/admanager/the-need-for-mobile-speed/}}, lo que subraya la importancia de optimizar los tiempos de respuesta. Asimismo, la actualización constante y relevancia de los datos es crucial; de lo contrario, la aplicación podría perder utilidad y credibilidad.

Por otro lado, la escalabilidad y la mantenibilidad son cualidades que garantizan la longevidad y adaptabilidad del software. Una aplicación escalable puede manejar un crecimiento en la demanda sin degradar su rendimiento, mientras que una aplicación mantenible facilita su reutilización, integración con sistemas existentes y actualización continua a lo largo de su ciclo de vida, reduciendo costos y esfuerzo en el proceso.

La arquitectura modular antes mencionada ofrece ventajas significativas para alcanzar estos objetivos. Por ejemplo, permite realizar modificaciones en una capa sin afectar a las demás. Esto significa que es posible actualizar una interfaz de usuario obsoleta en el \textit{frontend} sin alterar la lógica de negocio subyacente en el \textit{backend}, o incluso migrar a un sistema de base de datos más eficiente sin impactar la experiencia del usuario final.

A lo largo de los años, el desarrollo web se ha beneficiado de la aparición de \textit{frameworks} y librerías especializados, diseñados para simplificar la implementación, mejorar el rendimiento y facilitar el mantenimiento de las aplicaciones. Estas herramientas suelen estar orientadas no solo a una capa específica de la aplicación, sino también a nichos tecnológicos o tipos de aplicaciones particulares, como plataformas de comercio electrónico, redes sociales o sistemas de gestión de contenido.

En los siguientes apartados, se analizarán las diferentes tecnologías disponibles para el desarrollo web, comparando sus especializaciones, puntos fuertes y posibles limitaciones, con el objetivo de ofrecer una visión crítica que permita seleccionar las herramientas más adecuadas según los requisitos del proyecto más adelante.


\section{\textit{Frontend}}
La capa \textit{frontend} desempeña un papel fundamental al ser la responsable de la interfaz de usuario y la interacción del cliente con el sistema. Su principal objetivo es garantizar una experiencia de usuario óptima, centrada en la usabilidad, la accesibilidad y la fluidez en la interacción. Esta capa se construye sobre tres tecnologías esenciales que operan de manera coordinada: HTML (HyperText Markup Language), que define la estructura del contenido; CSS (Cascading Style Sheets), que establece el diseño visual y la presentación; y JavaScript, que implementa la interactividad y el dinamismo de la interfaz.

La elección de las tecnologías \textit{frontend} adecuadas es un aspecto crítico en el desarrollo de aplicaciones web, ya que influye directamente en el rendimiento, la escalabilidad y la mantenibilidad del proyecto. Según la encuesta de desarrolladores de Stack Overflow en 2024\footnote{\url{https://survey.stackoverflow.co/2024/technology}}, las tecnologías \textit{frontend} más utilizadas incluyen React, jQuery, Angular y Vue.js. Sin embargo, la selección de una u otra depende de múltiples factores, como las características del proyecto, las habilidades del equipo de desarrollo y los requisitos de escalabilidad. En este sentido, el estudio comparativo realizado por \cite{Sekhar_Emmanni_2023}, destaca la importancia de evaluar estas tecnologías en función de sus ventajas técnicas y su alineación con los objetivos del proyecto. 

En primer lugar, jQuery\footnote{\url{https://jquery.com/}} es una librería de JavaScript orientada a simplificar la manipulación del DOM, la gestión de eventos y las llamadas asíncronas. Su enfoque es transversal y de bajo nivel en comparación con otras soluciones más modernas, ya que no impone una arquitectura de aplicación ni un modelo de componentes. Aunque su uso ha disminuido en aplicaciones complejas, sigue siendo relevante en proyectos pequeños o usada en conjunto con \textit{frameworks}.

Por otro lado, React\footnote{\url{https://react.dev/}}, Angular\footnote{\url{https://angular.dev/}} y Vue.js\footnote{\url{https://vuejs.org/}} pertenecen a una generación posterior de herramientas centradas en la construcción de interfaces de usuario mediante componentes reutilizables y en la gestión del estado de la aplicación.

React, desarrollado por Meta, es una librería de JavaScript centrada en la construcción de interfaces de usuario mediante componentes. Su enfoque en la modularidad y el uso de un DOM virtual permite optimizar el rendimiento, reduciendo los tiempos de carga y mejorando la eficiencia, como se evidencia en el estudio de \cite{Sekhar_Emmanni_2023}. Su flexibilidad ha favorecido su adopción en aplicaciones complejas en la gestión del estado.

Angular, desarrollado y mantenido por Google, es un \textit{framework} completo basado en el patrón MVC (Modelo-Vista-Controlador). Aunque su rendimiento puede ser ligeramente inferior en comparación con otras tecnologías similares, su principal fortaleza radica en las herramientas integradas que ofrece para enrutamiento y gestión de formularios. Esta aproximación lo hace especialmente adecuado para aplicaciones empresariales de gran escala. 

Vue.js es un \textit{framework} de JavaScript que destaca por su simplicidad y facilidad de integración en proyectos existentes. Estas características junto con su suave curva de aprendizaje, la han convertido en una herramienta ideal para el desarrollo de prototipados y extensiones a aplicaciones existentes. 

La Tabla \ref{tab:frontend_technologies} clasifica estas tecnologías según su nivel de abstracción y su enfoque arquitectónico, diferenciando entre herramientas de manipulación directa del DOM y soluciones basadas en componentes.

\begin{table}[h!]
\centering
\small
\begin{tabularx}{\textwidth}{X X X p{4.5cm}}
\toprule
\textbf{Tecnología} & \textbf{Tipo} & \textbf{Nivel de abstracción / enfoque} & \textbf{Uso típico} \\
\midrule
jQuery & Librería & Manipulación directa del DOM, gestión de eventos y AJAX. No impone arquitectura ni modelo de componentes. & Proyectos pequeños, mantenimiento de código legacy, soporte puntual en aplicaciones existentes. \\
\midrule
React & Librería & Basado en componentes, DOM virtual, enfoque declarativo para la interfaz. & Aplicaciones complejas con gestión avanzada del estado y alta interactividad. \\
\midrule
Angular & \textit{Framework} completo & Arquitectura estructurada (basada en componentes y servicios), herramientas integradas (routing, formularios). & Aplicaciones empresariales de gran escala con necesidades estructurales claras. \\
\midrule
Vue.js & \textit{Framework} progresivo & Basado en componentes, integración gradual, enfoque flexible y ligero. & Prototipado, aplicaciones progresivas, integración en proyectos existentes. \\
\bottomrule
\end{tabularx}
\caption{Comparativa de tecnologías \textit{frontend}}
\label{tab:frontend_technologies}
\end{table}


\section{\textit{Backend}}

El \textit{backend} representa la columna vertebral de las aplicaciones web dinámicas, asumiendo la responsabilidad de implementar la lógica de negocio, procesar datos y facilitar la comunicación con servicios externos. A diferencia del \textit{frontend}, que se ejecuta en el navegador del usuario, el \textit{backend} opera en el servidor, donde gestiona operaciones críticas como la autenticación de usuarios, la validación y procesamiento de datos, la interacción con bases de datos y la integración con APIs de terceros. El diseño y la eficiencia de esta capa son factores determinantes para asegurar el rendimiento óptimo, la escalabilidad y la seguridad de las aplicaciones.

De acuerdo con los resultados de la encuesta de desarrolladores de Stack Overflow anteriormente mencionada, las tecnologías \textit{backend} más utilizadas en la actualidad incluyen Node.js, Express, Flask y Django. Cada una de estas tecnologías presenta características distintivas que las hacen adecuadas para diferentes tipos de proyectos. Además, como se detalla en el estudio comparativo de \textit{Frameworks} populares para \textit{backend} realizado por \cite{Buljic_Kadusic_Cvijanovic_Hadzajlic_Zivic_2025}, estos \textit{frameworksf} se basan en lenguajes de programación y tecnologías distintas, lo que influye en su rendimiento (medido en términos de latencia y throughput) y, por tanto, en su idoneidad para distintos escenarios de aplicación.

Node.js\footnote{\url{https://nodejs.org/es}} es un entorno de ejecución de JavaScript en el lado del servidor, diseñado específicamente para el desarrollo de aplicaciones escalables y de alto rendimiento. Como se muestra en los resultados del estudio comparativo realizado por \cite{Lei_Ma_Tan_2014}, su arquitectura basada en eventos y su modelo de entrada/salida (E/S) no bloqueante lo convierten en una opción ideal para aplicaciones que requieren manejar múltiples conexiones simultáneas, como sistemas en tiempo real o APIs. Además, Node.js cuenta con un ecosistema robusto, gracias a npm (Node Package Manager), que proporciona una amplia variedad de módulos para extender funcionalidades y agilizar el proceso de desarrollo.

Por su parte, Express\footnote{\url{https://expressjs.com/es/}} es un \textit{framework} minimalista para Node.js, diseñado para simplificar la creación de servidores web y APIs. Le permite a los desarrolladores construir aplicaciones robustas y eficientes con una configuración mínima, ofreciendo la flexibilidad necesaria para adaptarse a las necesidades específicas de cada proyecto.

Aunque Next.js\footnote{\url{https://nextjs.org/}} es reconocido principalmente como un \textit{framework} de React para el desarrollo de interfaces de usuario, su capacidad para implementar Server-Side Rendering (SSR) y rutas API lo posiciona como una herramienta híbrida que también puede gestionar lógica de \textit{backend}. Esta característica permite a los desarrolladores construir aplicaciones full-stack utilizando un único \textit{framework}, lo que simplifica la arquitectura del proyecto y reduce su complejidad.

En el ecosistema de Python, tanto Django\footnote{\url{https://www.djangoproject.com/}} como Flask\footnote{\url{https://flask.palletsprojects.com/en/stable/}} son \textit{frameworks} ampliamente utilizados para el desarrollo de aplicaciones web. Sin embargo, presentan diferencias significativas en cuanto a su enfoque y funcionalidades. Django es un \textit{framework} robusto y completo, que incluye una amplia gama de herramientas integradas, lo que lo hace especialmente adecuado para el desarrollo de sistemas empresariales complejos. En contraste, Flask se caracteriza por su ligereza y flexibilidad, siendo una opción preferida para proyectos más pequeños o aquellos que requieren una mayor personalización.

En la tabla \ref{tab:backend_technologies} se puede ver un resumen de la comparativa entre las diferentes tecnologías \textit{backend} que analizadas.

\begin{table}[h!]
\centering
\small
\begin{tabularx}{\textwidth}{X X X p{4.5cm} p{3cm}}
\toprule
\textbf{Framework} & \textbf{Lenguaje} & \textbf{Tipo} & \textbf{Características} & \textbf{Uso típico} \\
\midrule
Node.js & JavaScript & Entorno de ejecución & Arquitectura basada en eventos, E/S no bloqueante, escalable, ecosistema npm. & Aplicaciones en tiempo real, APIs. \\
\midrule
Express & JavaScript & \textit{Framework} minimalista & Minimalista, fácil configuración, ideal para APIs y servidores web. & APIs RESTful, aplicaciones web ligeras. \\
\midrule
Next.js & JavaScript / TypeScript & \textit{Framework} full-stack & SSR, rutas API, integración con React. & Aplicaciones full-stack, SSR. \\
\midrule
Django & Python & \textit{Framework} completo & Baterías incluidas, ORM, panel de administración, seguridad integrada. & Sistemas empresariales. \\
\midrule
Flask & Python & Micro-\textit{framework} & Ligero, flexible, fácil de extender. & Proyectos pequeños, APIs. \\
\bottomrule
\end{tabularx}
\caption{Comparativa de tecnologías \textit{backend}}
\label{tab:backend_technologies}
\end{table}

\section{Conclusión}

El cambio de paradigma en el desarrollo web actual está impulsado por la necesidad de crear aplicaciones más rápidas, interactivas y eficientes. En este contexto, surge una clara necesidad de adoptar tecnologías que no solo agilicen la creación de aplicaciones, sino que también permitan una fácil reutilización, integración y escalabilidad a largo plazo.

Por ello, los \textit{frameworks} y librerías modernos han demostrado ser herramientas esenciales para los desarrolladores. Permiten construir aplicaciones con una gran funcionalidad de manera rápida y eficiente, proporcionando herramientas listas para usar que reducen considerablemente el tiempo de desarrollo. Además, están altamente especializados, lo que significa que existen opciones adaptadas a prácticamente todas las necesidades, y basta con conocer las opciones, la aplicación y el entorno de desarrollo, para elegir una solución potente y adaptada. 